
\section{Una función, un fallo y una consecuencia}
La definición de fallo comienza por la función requerida y sus condiciones. Una señal fuera de un valor nominal no demuestra necesariamente pérdida de función; tampoco una pieza todavía capaz de funcionar satisface automáticamente todos los requisitos de seguridad o vida. Se distingue modo de fallo, causa y efecto. «Alta temperatura» puede ser un síntoma, mientras que pérdida de lubricación, degradación aerodinámica y error de sensor son hipótesis diferentes.

La consecuencia modifica la estrategia: un defecto evidente para la tripulación, una función oculta y una pérdida crítica no se tratan solo por su frecuencia. Detectar un fallo latente puede exigir una tarea funcional específica. El mantenimiento basado en fiabilidad examina si una tarea es técnicamente eficaz y qué consecuencia reduce; no consiste en aumentar todas las inspecciones por igual.
\section{Edad, condición y oportunidad de intervenir}
Un límite por edad puede ser eficaz si existe una relación pertinente entre edad y riesgo y una base que lo justifique. Si los fallos son esencialmente aleatorios, sustituir más frecuentemente puede no reducir la tasa del mismo modo e incluso introducir errores. Esto no autoriza a retirar tareas o límites obligatorios: evaluación técnica y vía de modificación son preguntas separadas.

El seguimiento de condición necesita una señal que cambie antes de perder función y un intervalo útil para detectar, decidir y actuar. Ese intervalo es incierto y depende de mecanismo y medida. Un sensor disponible no garantiza capacidad predictiva: se necesita conocer error, condiciones y significado de la alarma.
\section{Mantenibilidad y tiempos de indisponibilidad}
Accesos, modularidad, diagnóstico y útiles influyen en tiempos y errores de restauración. Reparación activa no incluye necesariamente espera por repuesto, aprobación o transporte. Al comparar disponibilidad se declara qué estados y tiempos se cuentan, como ilustra P30. Una mejora puede proceder de menos fallos, menor tiempo técnico o menor espera; son acciones diferentes. Un indicador agregado puede ocultar defectos críticos repetidos si la mayoría de la flota acumula muchas horas sin incidentes.
\repaso{¿Toda tarea preventiva reduce necesariamente riesgo?}{Debe detectar o evitar el modo relevante con eficacia y ejecutarse correctamente. También puede introducir daño o errores; eso se valora junto a sus beneficios y requisitos obligatorios.}
